\section{Conclusions}
\label{sec:conclusions}

This study examines the energetic response of a lid-driven cavity to superelliptic bottom-wall shaping. For $n=2\to4$ at $\alpha=0$, fluid area increases by $7.93\%$, but integrated dissipation rises by only $0.0218\%$, $0.3249\%$, and $0.5909\%$ at $Re=100$, $400$, and $1000$. An area-only reference overpredicts these contrasts by factors of $364$, $24$, and $13$. Dissipation per unit fluid area decreases by about $7\%$, showing that area change alone cannot explain the response.

An exact spatial partition separates dissipation in the newly opened shoulder region from redistribution over the region common to both geometries. This $\mathrm{d}A/\mathrm{d}B$ decomposition, which constitutes the central original contribution of the present work, provides direct physical insight into the mechanism underlying the sub-extensive response: at $Re=400$, the added region occupies $7.4\%$ of the fluid area but contributes only $0.2\%$ of the reference dissipation, while the common-region redistribution is itself a near-cancellation of opposing lid-band and core-band contributions. The common-region contribution is negative at $Re=100$ and positive at $Re=400$ and $1000$, reflecting a Reynolds-dependent balance between near-wall and core changes. The net contrast remains positive throughout. The grid robustness of this mechanism is assessed in Table~SM.XV and Fig.~SM.5: all four bands are recomputed with the same exact-partition code at every resolution $N=128$--$512$; the cancellation index $C=\sum|\text{bands}|/|\mathrm{d}A|$ grows from $4.0$ at $Re=100$ to $14.3$ at $Re=1000$, but the net contrast $\Delta D=\mathrm{d}A+\mathrm{d}B$ remains stable to within $2.4\%$ across all grids because the grid-induced changes in $\mathrm{d}A$ and $\mathrm{d}B$ partially cancel. The sign and order of magnitude of $\mathrm{d}A$ and $\mathrm{d}B$ are therefore robust on every grid; their individual magnitudes carry a grid sensitivity that rises with $Re$ (the $\mathrm{d}A$ span is $1.9\%$, $8.2\%$ and $27\%$ of its value at $Re=100$, $400$ and $1000$). At $Re=100$, the energy-balance defect varies by $3.8$ points between $N=128$ and $N=384$ while $I_D(2,4)$ remains within $3.4\times10^{-4}$~pp (Table~SM.XIX), confirming that the defect is a common-mode error that cancels in the paired contrast.

The response is concentrated in the $n=2\to4$ transition at all Reynolds numbers studied, which carries about $96.2\%$ of the summed absolute contrasts at $Re=400$. The full exponent sequence is now available at all three Reynolds numbers: at $\alpha=0$ the $n=4\to6$ contrast is $+0.00296\%$, $+0.0129\%$, and $+0.0103\%$, and the $n=6\to8$ contrast is $+0.00042\%$, $+6\times10^{-5}\%$, and $-0.0037\%$ at $Re=100$, $400$, and $1000$. Resolution is judged transition by transition against each contrast's own $N=128$--$384$ span. At $Re=1000$ both trailing steps are resolved, the $(6,8)$ step by a factor of $34$ and with the same negative sign on all three grids, so it is reported as the first resolved negative trailing step. At $Re=400$ that step lies inside its own span, so its sign is undetermined; at $Re=100$ the $(4,6)$ step cannot be resolution-tested, because its two coarser members mix Mach regimes, and it too is quoted as an upper bound. At $\alpha=0.1$ the two trailing transitions are negative and individually resolved at both $Re=400$ and $Re=1000$, so wall undulation fixes their sign already at moderate Reynolds number. Increasing undulation amplitude also weakens the primary contrast. Across $N=128$, $256$, and $384$, the primary $\alpha=0$ contrast varies by at most $2.3\%$ of its own value, although absolute integrals retain about $1\%$ grid sensitivity. Independent Taylor--Hood finite-element calculations reproduce the primary contrast to $0.002$ percentage points at $Re=400$ and to $0.008$ percentage points at $Re=1000$, on the finest meshes of each solver. Scope is limited to the tested superellipse family and single undulation count; the unmatched Mach numbers of the coarse $n\ge6$ runs at $Re=100$, together with residual grid sensitivity, preclude universal scaling or grid-independent claims. Numerical uncertainties on $I_D(2,4)$ are $3.4\times10^{-4}$, $2.4\times10^{-3}$, and $1.4\times10^{-2}$ percentage points at $Re=100$, $400$, and $1000$; the trailing contrasts have their own, generally smaller spans, collected in Table~SM.X, and are $6.9\times10^{-4}$ and $1.1\times10^{-4}$ percentage points at $Re=1000$, well below the corresponding magnitudes.
